\documentclass[11pt,a4paper]{article}
\usepackage[margin=2.15cm]{geometry}
\usepackage[T1]{fontenc}
\usepackage{lmodern}
\usepackage{amsmath,amssymb,mathtools}
\usepackage{enumitem}
\usepackage{fancyhdr}
\usepackage{booktabs}
\usepackage{tikz}
\usepackage[hidelinks]{hyperref}
\usepackage{microtype}

\setlength{\parindent}{0pt}
\setlength{\parskip}{0.4em}
\setlist[enumerate]{leftmargin=*,itemsep=0.3em,topsep=0.3em}
\pagestyle{fancy}
\setlength{\headheight}{14pt}
\fancyhf{}
\lhead{Physics Practice Examination --- Answer Key}
\rhead{\thepage}
\renewcommand{\headrulewidth}{0.4pt}

\newcommand{\problem}[2]{\subsection*{Problem #1 \; \textnormal{\small[#2]}}}
\newcommand{\unit}[1]{\,\mathrm{#1}}
\newenvironment{parts}{\begin{enumerate}[label=\textbf{\alph*)}]}{\end{enumerate}}

\title{\vspace{-1.4cm}\bfseries Physics Practice Examination\\Answer Key and Worked Solutions}
\author{}
\date{}

\begin{document}
\maketitle
\vspace{-2em}

%======================================================================
\problem{1}{FY1}
Theory gives $T=2\pi\sqrt{(l+r)/g}$, hence
\[
  T^2 = \frac{4\pi^2}{g}\,l + \frac{4\pi^2 r}{g},
\]
a straight line in $l$ with slope $k=4\pi^2/g$ and intercept
$b=4\pi^2r/g$. It must be $T^2$ that is plotted: $T$ against $l$ is a
square-root curve, and fitting it with a straight line would be a modelling
error. The intercept is nonzero precisely because of the bob radius, so the
model must keep $r$.

\textbf{(a)} As above: $k=4\pi^2/g$, $b=4\pi^2r/g$.

\textbf{(b)} Least squares (unweighted; no uncertainties are quoted):
\[
  T^2 = 4.0805\,l + 0.02242\qquad(T^2\ \text{in}\ \unit{s^2},\ l\ \text{in m}),
\]
so
\[
  g=\frac{4\pi^2}{4.0805\unit{s^2/m}} = 9.67\unit{m/s^2}.
\]
Residuals in $T^2$ are $-3.1,\,1.7,\,3.6,\,0.1,\,-2.3\unit{ms^2}$, giving
$s_y=3.2\cdot10^{-3}\unit{s^2}$ with $\nu=3$ degrees of freedom and
\[
  s_k=\frac{s_y}{\sqrt{\sum (l_i-\bar l)^2}} = 5.1\cdot10^{-3}\unit{s^2/m},
  \qquad \frac{s_k}{k}=0.12\%.
\]
Since $g\propto k^{-1}$, $\Delta g/g=\Delta k/k$, so
\[
  \boxed{g = (9.67\pm0.01)\unit{m/s^2}}\qquad(0.12\%).
\]
This is a \emph{statistical} uncertainty: it describes the scatter of these
five points only.

\textbf{(c)} Intercept $b=(4\pi^2/g)r$, with standard uncertainty
$s_b=s_y\sqrt{1/n+\bar l^2/\sum(l_i-\bar l)^2}=3.4\cdot10^{-3}\unit{s^2}$.
Then
\[
  r=\frac{b\,g}{4\pi^2}=5.5\unit{mm},\qquad
  \Delta r\approx0.8\unit{mm},
\]
and $b/s_b=6.6$. The intercept is significantly positive (about
$6\sigma$), consistent with a bob of roughly $1\unit{cm}$ diameter.
Fitting a model with $r=0$ imposed would be a modelling error.

\textbf{(d)} The finite-amplitude correction is
\[
  T=T_0\Bigl(1+\frac{\theta_0^2}{16}\Bigr),\qquad
  \theta_0=20.0^\circ=0.3491\unit{rad},\quad
  \frac{\theta_0^2}{16}=0.00762 .
\]
Since $g\propto1/T^2$, the fitted slope is too large by
$(1+0.00762)^2=1.01529$, so
\[
  g_{\text{wrong}}=\frac{9.81}{1.01529}=9.663\unit{m/s^2},
\]
i.e.\ $1.50\%$ low, against the observed $1.4\%$. The correction therefore
accounts for the discrepancy, and applying it gives
\[
  g_{\text{corr}} = 9.67\cdot1.01529 = \boxed{9.82\unit{m/s^2}} .
\]
This is a \emph{systematic} error: it shifts $g$ in one direction, it is
\emph{not} reduced by more repetitions, and it is invisible in the fit
residuals --- the points lie on a perfect straight line with the wrong
slope. The statistical uncertainty from part (b) therefore does not cover
the deviation, which is why the result looked inconsistent. Removing it:
\begin{itemize}
\item Vary the amplitude ($5^\circ$, $10^\circ$, $20^\circ$) and extrapolate
      $T$ to $\theta_0\to0$. This is the correct procedure.
\item Work at such small amplitudes that the correction is below the
      timing resolution.
\item With the photogate, time \emph{many} periods and divide.
\end{itemize}
The reading uncertainty $\pm0.5^\circ$ alone gives only about
$\pm0.1\%$ in $g$ (differentiate $(1+\theta_0^2/16)^2$), so this is not a
reading-uncertainty problem at all: it is a model error.

%======================================================================
\problem{2}{FY2}
\textbf{(a)} The kinetic energy of the air of mass $m$ passing the disc in
time $\Delta t$ is $\tfrac12mv^2$ with $m=\rho Av\Delta t$; hence
$P=\tfrac12\rho Av^2\cdot v=\tfrac12\rho A v^3$. Dimensional check:
$\unit{kg/m^3}\cdot\unit{m^2}\cdot\unit{m^3/s^3}=\unit{J/s}$.
With $A=\pi D^2/4=1.131\cdot10^{4}\unit{m^2}$ and $v=12\unit{m/s}$,
\[
  P_w = 1.17\cdot10^{7}\unit{W} = \boxed{11.7\unit{MW}}.
\]
Air cannot be stopped dead: mass conservation forces it to leave the rotor
with some speed, so the extracted power is
$P=\tfrac12\dot m(v_1^2-v_2^2)$ with $v_2>0$. Optimising over $v_2$ gives
the Betz limit $C_{p,\max}=16/27=0.593$. At $12\unit{m/s}$ the turbine
delivers $2.50\unit{MW}$, i.e.\ $21\%$ of $P_w$ and $36\%$ of the Betz
limit; it is well below Betz because $12\unit{m/s}$ is already at the rated
point, where the controller deliberately spills power.

\textbf{(b)} With $\bar v=7.5\unit{m/s}$,
$\tfrac12\rho A\bar v^3=2.863\unit{MW}$, and
\[
  \overline{v^3}=\frac6\pi\bar v^3 = 1.9099\,\bar v^3
  \;\Longrightarrow\;
  \overline{P}_{\text{avail}}=\tfrac12\rho A\overline{v^3}=5.468\unit{MW}.
\]
Hence
\[
  C_p = \frac{2.50\unit{MW}}{5.468\unit{MW}} = \boxed{0.457} < 0.593 .
\]
Using $\bar v^3$ would give $C_p=0.87>16/27$, which is impossible. Because
$v\mapsto v^3$ is convex, $\overline{v^3}\ge\bar v^3$, so substituting the
mean speed \emph{underestimates} the available power and
\emph{overestimates} $C_p$. The two coincide only for a perfectly steady
wind.

\textbf{(c)} $E=P_{\text{rated}}\,t\,\text{CF}
=2.50\unit{MW}\cdot8760\unit{h}\cdot0.35=7.67\unit{GWh}$, so
\[
  N=\frac{0.010\cdot90\unit{TWh}}{7.67\unit{GWh}} \approx
  \boxed{1.2\cdot10^{2}\ \text{turbines}} .
\]
CO$_2$ avoided: $7.665\cdot10^{6}\unit{kWh}\cdot0.90\unit{kg/kWh}
=6.9\cdot10^{6}\unit{kg} \approx \boxed{6.9\cdot10^{3}\unit{t}}$ per year.

\textbf{(d)} $\overline{v^3}=\tfrac12(5^3+10^3)=562.5$ whereas
$\bar v^3=7.5^3=421.9$, so the alternating site yields
\[
  \frac{562.5}{421.875}=\boxed{\frac43}
\]
times more power --- a $33\%$ gain from the same mean speed. This is why
the average power involves $\overline{v^3}$. A real turbine does not keep
producing more at ever higher wind speed because it is rated at a design
wind speed ($\approx12$--$13\unit{m/s}$ here): above it, pitch control sheds
power to protect the structure, and at storm speeds the machine is shut
down.

%======================================================================
\problem{3}{FY3}
\textbf{(a)} $\rho_0=p_0M/(RT_0)
=101.3\cdot10^{3}\cdot0.0290/(8.314\cdot293.15)
=\boxed{1.205\unit{kg/m^3}}$.

\textbf{(b)} Buoyancy must balance the total weight,
\[
  V\rho_0 g = V\rho_i g + (m_{\text{env}}+m_{\text{bask}}+m_{\text{pay}})g,
\]
so the required internal density is
\[
  \rho_i=\rho_0-\frac{m_{\text{tot}}}{V}
  =1.205-\frac{180+90+350}{2.2\cdot10^3}
  =1.205-0.2818=0.9235\unit{kg/m^3}.
\]
At constant pressure $\rho\propto1/T$, hence
\[
  T_i=T_0\frac{\rho_0}{\rho_i}=293.15\cdot\frac{1.2053}{0.9235}
  =382.6\unit{K}=\boxed{109^\circ\text{C}}.
\]

\textbf{(c)} At $T_i=393.15\unit{K}$,
$\rho_i=\rho_0T_0/T_i=0.8991\unit{kg/m^3}$, so the total mass that can be
carried is
\[
  m_{\text{tot}}=V(\rho_0-\rho_i)=2.2\cdot10^{3}\cdot0.3063=674\unit{kg},
\]
leaving a payload of $674-270=\boxed{4.0\cdot10^{2}\unit{kg}}$.

\textbf{(d)} The burner must supply the enthalpy to raise the entering air
from $T_0$ to $T_i$:
\[
  \dot m=\rho_i\dot V=0.9235\cdot0.45=0.4156\unit{kg/s},
\]
\[
  P=\dot m\,c_p\,(T_i-T_0)=0.4156\cdot1005\cdot(382.6-293.15)
  =\boxed{3.7\cdot10^{4}\unit{W}\approx37\unit{kW}}.
\]
Approximations: (i) constant pressure, so $c_p$ is the right specific
heat; (ii) the air leaves at the same temperature at which it enters, i.e.\
no heat is lost through the fabric and the mixing/leakage losses are
ignored. Every one of these makes the estimate a \emph{lower bound}: real
envelopes lose heat and the expelled air carries away more enthalpy than
assumed. That $37\unit{kW}$ is unrealistically small for a
$2200\unit{m^3}$ balloon (typical burners are hundreds of kW) is exactly
the size of the neglected losses.

%======================================================================
\problem{4}{FY4}
\textbf{(a)} With B descending and A moving up the plane, the driving force
is
\[
  m_Bg-m_Ag\sin\theta = 29.43-19.62=9.81\unit{N},
\]
while the maximum static friction is
$\mu_s m_Ag\cos\theta=0.25\cdot33.98=8.50\unit{N}$. Since $9.81>8.50$, A
does move as assumed. Then $f_k=\mu_k m_Ag\cos\theta=6.80\unit{N}$ and
\[
  a=\frac{m_Bg-m_Ag\sin\theta-f_k}{m_A+m_B}
  =\frac{9.81-6.80}{7.00}=\boxed{0.430\unit{m/s^2}},
\]
\[
  T=m_B(g-a)=3.00(9.81-0.430)=\boxed{28.1\unit{N}}.
\]
Check on A: $T-m_Ag\sin\theta-f_k=28.14-19.62-6.80=1.72\unit{N}$ and
$m_Aa=1.72\unit{N}$. \checkmark

\textbf{(b)} $v=\sqrt{2ah}=\sqrt{2\cdot0.4305\cdot0.80}
=\boxed{0.830\unit{m/s}}$.
Energy balance (all terms in J):
\[
  \underbrace{m_Bgh}_{23.54}
  =\underbrace{m_Agh\sin\theta}_{\Delta U_A=15.70}
  +\underbrace{f_kh}_{W_f=5.44}
  +\underbrace{\tfrac12(m_A+m_B)v^2}_{K=2.41},
\]
and $15.70+5.44+2.41=23.54$. \checkmark

\textbf{(c)} After B lands, only gravity and friction act on A:
\[
  a_2=g(\sin\theta+\mu_k\cos\theta)=9.81(0.500+0.1732)
  =6.60\unit{m/s^2}
\]
(deceleration), so
\[
  s=\frac{v^2}{2a_2}=\frac{0.6887}{13.21}=\boxed{5.2\unit{cm}}.
\]
A stays put only if $\mu_s\ge\tan\theta=0.577$; here $\mu_s=0.25$, so A
slides back down.

\textbf{(d)} Sliding down, $a_3=g(\sin\theta-\mu_k\cos\theta)
=9.81(0.500-0.1732)=3.21\unit{m/s^2}$, and the distance back to the point
where B landed is $s=5.2\unit{cm}$. So A's speed there is
\[
  v_A=\sqrt{2a_3s}=\sqrt{2\cdot3.206\cdot0.05215}
  =\boxed{0.578\unit{m/s}} .
\]
At that instant the string snaps taut. The jerk is so brief that the
impulses of gravity and friction are negligible; the tension impulse
$J=\int T\,dt$ is internal to the pair A+B, so momentum is conserved. B is
at rest on the floor at that moment and is jerked into motion, so the two
blocks end up with a common speed
\[
  (m_A+m_B)V=m_Av_A \;\Longrightarrow\;
  V=\frac{m_A}{m_A+m_B}v_A=\frac47\cdot0.578=\boxed{0.33\unit{m/s}},
\]
with an impulse $J=m_BV=\boxed{0.99\unit{N\,s}}$ transferred from A to B
(the same $J$ equals $m_A(v_A-V)$). The kinetic energy is
\[
  K_{\text{before}}=\tfrac12m_Av_A^2=0.669\unit{J},\qquad
  K_{\text{after}}=\tfrac12(m_A+m_B)V^2=0.382\unit{J},
\]
so $0.287\unit{J}$ is lost in the inelastic jerk, i.e.\ the fraction
\[
  \frac{|\Delta K|}{K}=\frac{m_B}{m_A+m_B}=\frac37=0.43,
\]
exactly as in a perfectly inelastic collision, and this fraction does not
depend on the masses' arrangement.

What happens next is a nice check on the sign conventions: A is moving
down the plane and B upward, so gravity now opposes the motion, and with
the string taut
\[
  a=\frac{m_Ag\sin\theta-m_Bg-\mu_km_Ag\cos\theta}{m_A+m_B}
  =\frac{19.62-29.43-6.80}{7.00}=-2.4\unit{m/s^2},
\]
i.e.\ the pair decelerates. Starting from $V=0.33\unit{m/s}$ it stops
after only $0.023\unit{m}$ and $0.14\unit{s}$, whereupon B falls back onto
the floor and the string goes slack again. The blocks therefore descend in
a sequence of small jerks, losing energy at every one; the total energy lost
is not recoverable, which is why a taut inelastic string is not a
conservative constraint.

%======================================================================
\problem{5}{FY5}
\textbf{(a)} $r=R_E+h=6.771\cdot10^{6}\unit{m}$ and
$GM_E=3.984\cdot10^{14}\unit{m^3/s^2}$.
\[
  v=\sqrt{\frac{GM_E}{r}}=\sqrt{\frac{3.984\cdot10^{14}}{6.771\cdot10^{6}}}
  =\boxed{7.67\unit{km/s}},
\]
\[
  T=\frac{2\pi r}{v}=5.55\cdot10^{3}\unit{s}=\boxed{92.4\unit{min}},
\]
\[
  g(r)=\frac{GM_E}{r^2}=\boxed{8.69\unit{m/s^2}}.
\]

\textbf{(b)} Circular motion requires $mv^2/r=GM_Em/r^2$, hence
$K=\tfrac12mv^2=\tfrac12 GM_Em/r$. Since $U=-GM_Em/r$, we get
$K=-\tfrac12U$ and
\[
  E=K+U=\tfrac12U=-\frac{GM_Em}{2r}.
\]
With $m=4.2\cdot10^{5}\unit{kg}$,
\[
  E=-\frac{3.984\cdot10^{14}\cdot4.2\cdot10^{5}}{2\cdot6.771\cdot10^{6}}
  =-1.24\cdot10^{13}\unit{J}
  =\boxed{-12.4\unit{TJ}}.
\]
Escape speed: $v_{\text{esc}}=\sqrt{2GM_E/r}=10.8\unit{km/s}=\sqrt2\,v$.
The two numbers differ by exactly this factor because escaping means
supplying $|E|=K$: you need to raise $E$ from $-K$ to $0$.

\textbf{(c)} $E\propto-1/r$, so with $\Delta r=-1.00\unit{km}$,
\[
  \Delta E \approx -\frac{|E|\,\Delta r}{r}
  =-\frac{1.236\cdot10^{13}\cdot(-10^{3})}{6.771\cdot10^{6}}
  = \boxed{+1.8\unit{GJ}} .
\]
The total energy \emph{decreases} by $1.8\unit{GJ}$: that energy must be
removed from the orbit (by drag or by a retro burn). Of it, the potential
energy falls by $2\times1.8\unit{GJ}=3.7\unit{GJ}$ while the kinetic energy
rises by $1.8\unit{GJ}$ (since $K=-E$). From $v=\sqrt{GM/r}$,
$\mathrm dv=-\tfrac12(v/r)\mathrm dr$, so
\[
  \Delta v = +\tfrac12\frac{v}{r}\cdot10^{3}=\boxed{+0.57\unit{m/s}},
\]
and $T\propto r^{3/2}$ gives $\Delta T=\tfrac32(T/r)\Delta r=-1.2\unit{s}$.
So the station speeds \emph{up} by $0.57\unit{m/s}$ while losing energy: a
retro burn lowers the orbit and converts more potential energy into kinetic
energy than the burn removes. A prograde burn instead raises the orbit and
slows the station. This is the reverse of the intuition from driving a car.

\textbf{(d)} $a=\omega^2R_s=0.256^2\cdot50.0=3.28\unit{m/s^2}$. At the
head, $R=50.0-1.80=48.2\unit{m}$, so
\[
  a_{\text{head}}=\omega^2\cdot48.2=3.16\unit{m/s^2},
\]
i.e.\ $1-48.2/50.0=3.6\%$ weaker. The relative gradient is
$\Delta a/a=\Delta R/R_s$, so to keep it below a few per cent $R_s$ must be
large compared with human height: a healthy habitat needs a radius of
tens or hundreds of metres. Small habitats suffer Coriolis artefacts,
gradients and vertigo.

\textbf{(e)} In the habitat frame the pendulum's equilibrium direction is
radial and the restoring "gravity" is the centrifugal acceleration at its
own location; taking the bob's radius as $\approx50\unit{m}$, so that
$g_{\text{eff}}=a=3.28\unit{m/s^2}$,
\[
  T=2\pi\sqrt{\frac{L}{a}}=2\pi\sqrt{\frac{3.00}{3.28}}
  =\boxed{6.0\unit{s}},
\]
against $2\pi\sqrt{3.00/9.81}=3.5\unit{s}$ on Earth: longer by
$\sqrt{9.81/3.28}=\sqrt3$, since the period scales as $g^{-1/2}$. (Strictly
the bob sits $3.00\unit{m}$ inboard of the feet, where $a$ is $6\%$ smaller;
using $R=47$ and $50\unit{m}$ gives $6.2$ and $6.0\unit{s}$ respectively.
The point of the question is the scaling, not the third digit.)

%======================================================================
\problem{6}{FY5}
\textbf{(a)} $s=A\sin\omega t$ with
$\omega=2\pi f=2.765\cdot10^{3}\unit{rad/s}$ and $A=5.00\cdot10^{-4}\unit{m}$:
\[
  v_{\max}=\omega A=\boxed{1.38\unit{m/s}},\qquad
  a_{\max}=\omega^2A=\boxed{3.82\cdot10^{3}\unit{m/s^2}}\approx390\,g.
\]
The air molecules oscillate about fixed positions; they do not travel with
the wave. For a pipe closed at one end, resonances occur at
$l_n+\Delta l=(2n-1)\lambda/4$, so
\[
  l_2-l_1=\frac{\lambda}{2}
  \;\Longrightarrow\;
  \lambda=2(0.572-0.182)=0.780\unit{m},
\]
\[
  v=f\lambda=440\cdot0.780=\boxed{343\unit{m/s}} .
\]

\textbf{(b)} The ideal quarter-wave length is $\lambda/4=0.195\unit{m}$, but
the resonance is at $0.182\unit{m}$: the air just outside the open end also
moves, effectively lengthening the pipe. In this convention
$l_{\text{resonance}}+\Delta l=\lambda/4$, so
\[
  \Delta l=\frac{\lambda}{4}-l_1=0.195-0.182=1.3\unit{cm}.
\]
With $\Delta l=0.6r$, $r=1.3/0.6=2.2\unit{cm}$ (diameter about
$4.3\unit{cm}$) --- a plausible resonance tube. $\Delta l$ cancels in
$l_2-l_1$ because it is the same for every resonance, but it does not
cancel in $\Delta l$ itself, which is why the absolute positions are what
determine the tube geometry.

\textbf{(c)} $f_2=f\pm4.0\unit{Hz}$, i.e.\ $436$ or $444\unit{Hz}$.
To decide: load the unknown fork slightly (a dab of plasticine). Loading
lowers its frequency, so the beat frequency increases if $f_2$ was already
above $440\unit{Hz}$, and decreases if it was below. Beats occur at
$|f_1-f_2|=4.0$ per second, so in $5.0\unit{s}$ there are
$\boxed{20}$ beats.

\textbf{(d)} With the same end correction,
$\lambda_{\text{CO}_2}=4(l_1+\Delta l)=4(0.140+0.013)=0.612\unit{m}$, so
\[
  v_{\text{CO}_2}=440\cdot0.612=\boxed{269\unit{m/s}}.
\]
Kinetic theory gives $v=\sqrt{\gamma RT/M}$:
\[
  v_{\text{air}}=\sqrt{\frac{1.40\cdot8.314\cdot293}{0.0290}}
  =343\unit{m/s},\qquad
  v_{\text{CO}_2}=\sqrt{\frac{1.30\cdot8.314\cdot293}{0.0440}}
  =268\unit{m/s}.
\]
The measured value agrees within $0.4\%$. The speed drops because the
larger molar mass dominates the slightly smaller $\gamma$: $v\propto
\sqrt{\gamma/M}$, and $M$ rises by a factor $1.52$ while $\gamma$ falls by
only $7\%$. The beat frequency of part (c) is a property of the two forks,
not of the medium, so it is unchanged; only the wavelength changes.

%======================================================================
\problem{7}{FY6}
\textbf{(a)} $R_1\parallel R_2$: $R_p=18\cdot22/40=9.90\unit{\Omega}$.
Total $R_{\text{tot}}=9.90+15.0+3.0=27.9\unit{\Omega}$, so
\[
  I=\frac{12.0}{27.9}=\boxed{0.430\unit{A}}.
\]
Terminal voltage $U=E-Ir=12.0-1.29=\boxed{10.7\unit{V}}$. The parallel
voltage is $U_p=IR_p=4.258\unit{V}$, hence
\[
  I_{R_1}=\frac{4.258}{18}=\boxed{0.237\unit{A}},\qquad
  I_{R_2}=\frac{4.258}{22}=\boxed{0.194\unit{A}},
\]
and $0.237+0.194=0.430\unit{A}$. \checkmark\ The capacitor carries $U_p$:
\[
  Q=CU_p=2.2\cdot10^{-6}\cdot4.258=\boxed{9.4\unit{\mu C}},\qquad
  W=\tfrac12CU_p^2=\boxed{2.0\cdot10^{-5}\unit{J}}.
\]

\textbf{(b)} Powers:
\[
\begin{array}{ll}
R_3: I^2R_3=2.77\unit{W} & R_1:U_p^2/R_1=1.01\unit{W}\\
R_2:U_p^2/R_2=0.824\unit{W} & r: I^2r=0.555\unit{W}
\end{array}
\]
summing to $5.16\unit{W}$, while $EI=12.0\cdot0.430=5.16\unit{W}$.
\checkmark\ Efficiency (external resistors only):
\[
  \eta=\frac{2.77+1.01+0.824}{5.16}=0.892=\boxed{89\%}.
\]

\textbf{(c)} With the voltmeter across the parallel pair, $R_1$ is loaded:
\[
  R_1'=\frac{R_1R_V}{R_1+R_V},\quad
  R_p'=\frac{R_1'R_2}{R_1'+R_2},\quad
  I'=\frac{E}{R_p'+R_3+r},\quad U'=I'R_p' .
\]
As $R_V\to\infty$, $U=4.258\unit{V}$.
$R_V=200\unit{\Omega}$: $R_1'=16.51$, $R_p'=9.45\unit{\Omega}$,
$I'=0.437\unit{A}$, $U'=\boxed{4.13\unit{V}}$ (3\% low).
$R_V=30\unit{\Omega}$: $R_1'=11.25$, $R_p'=7.44\unit{\Omega}$,
$I'=0.471\unit{A}$, $U'=\boxed{3.51\unit{V}}$ (18\% low).
The reading is always too small because the voltmeter draws current,
lowering the equivalent resistance of the pair and therefore its voltage
drop. Solving $U'(R_V)=0.99\,U$ gives
$R_V\approx6.3\cdot10^{2}\unit{\Omega}$, i.e.\ some $60\times$ the
$9.9\unit{\Omega}$ it is in parallel with. The practical rule is
$R_V\gg R_{\text{measured}}$; real meters have $\unit{M\Omega}$ inputs, so
the error vanishes.

\textbf{(d)} A diode's $I$--$U$ characteristic is exponential: a small
voltage change produces a huge current change, so direct connection to
$9.0\unit{V}$ would destroy the LED. The series resistor sets the current:
\[
  R=\frac{U-U_f}{I}=\frac{9.0-2.0}{0.020}=\boxed{350\unit{\Omega}},
  \qquad P_R=I^2R=\boxed{0.14\unit{W}} .
\]
Reversed, the diode is in the blocking direction: essentially no current
flows (only $\unit{nA}$ leakage), the LED does not light, and nothing is
damaged. As the battery discharges its terminal voltage falls, so
$U_R=U_{\text{batt}}-U_f$ falls, the current $I=U_R/R$ falls and the LED
dims. Because $U_f$ is nearly pinned by the exponential characteristic,
almost the whole voltage change appears across $R$, so the dimming is
gradual.

%======================================================================
\problem{8}{FY7}
\textbf{(a)} $f=3000/60=50\unit{Hz}$, $\omega=2\pi f=314\unit{rad/s}$.
With the coil normal at angle $\theta=\omega t$ to $\mathbf B$,
\[
  \Phi=BA\cos\omega t,\qquad
  e=-N\frac{d\Phi}{dt}=NBA\omega\sin\omega t .
\]
\[
  \hat e=NBA\omega
  =100\cdot0.85\cdot1.22\cdot10^{-2}\cdot314
  =\boxed{3.3\cdot10^2\unit{V}},
\]
\[
  e_{\text{rms}}=\frac{\hat e}{\sqrt2}=\boxed{2.3\cdot10^2\unit{V}}.
\]
Over a quarter turn, $\bar e=N\Delta\Phi/\Delta t$ with $\Delta\Phi=BA$
and $\Delta t=T/4=1/(4f)$, giving
\[
  \bar e = 4NBAf = \frac{2}{\pi}\hat e .
\]
Hence $\hat e/\bar e=\pi/2=1.57$. The reason is that $e(t)$ is sinusoidal:
a quarter turn covers the quarter of the sine from zero to its peak, whose
mean is $(2/\pi)$ times the peak, just as for any sinusoid. (The mean of
$|e|$ over a full period is also $(2/\pi)\hat e$, while the rms is
$\hat e/\sqrt2$; these three averages must not be confused.)

\textbf{(b)} Without a transformer, at $U=230\unit{V}$,
\[
  I=\frac{P}{U}=\frac{50\cdot10^{6}}{230}=2.2\cdot10^{5}\unit{A},\qquad
  P_{\text{loss}}=I^2R=\boxed{4.7\cdot10^{11}\unit{W}},
\]
which is nearly $10^{4}$ times the generated power --- utterly impossible.
That is itself the answer: without a transformer the line cannot deliver
the power at all. In general
\[
  P_{\text{loss}}=I^2R=\frac{P^2R}{U^2},
\]
so the loss scales as $U^{-2}$.
With $U'=110\unit{kV}$: $I=455\unit{A}$ and
\[
  P_{\text{loss}}=(455)^2\cdot10=\boxed{2.1\cdot10^{6}\unit{W}
  =2.1\unit{MW}},
\]
which is $\boxed{4.1\%}$ of $50\unit{MW}$. A transformer needs a
\emph{changing} flux: for DC, $\Phi\propto I$ is constant, so the secondary
emf $N_2\,d\Phi/dt=0$. DC cannot be transformed, which is why HVDC links
convert to AC or use power electronics at both ends. Raising the voltage
is a cost trade: line loss falls as $U^{-2}$ while insulation, tower height
and clearances cost grows with $U$. The optimum balances the two. It is not
about transformer efficiency, which already exceeds $99\%$.

%======================================================================
\problem{9}{FY7}
\textbf{(a)} Let the rod slide down the incline at speed $v$. The field is
perpendicular to the rail plane, so
$\mathcal{E}=BLv$, $I=BLv/R$, and the magnetic force on the rod is
$F=BIL=B^2L^2v/R$, directed \emph{up} the incline.

\emph{i.} Terminal speed follows from $mg\sin\theta=B^2L^2v_t/R$:
\[
  v_t=\frac{mgR\sin\theta}{B^2L^2}
  =\frac{0.20\cdot9.81\cdot0.30\cdot0.500}{1.2^2\cdot0.30^2}
  =\boxed{2.27\unit{m/s}} .
\]
\emph{ii.}
$I=BLv_t/R=1.2\cdot0.30\cdot2.27/0.30=\boxed{2.73\unit{A}}$. Electrical
power
\[
  P_e=I^2R=2.73^2\cdot0.30=\boxed{2.23\unit{W}},
\]
while the mechanical power released is
$mg\sin\theta\,v_t=0.981\cdot2.27=2.23\unit{W}$. \checkmark

\emph{iii.} Pushed up at $2.0\unit{m/s}$, the motional emf is
$\mathcal{E}=BLv=1.2\cdot0.30\cdot2.0=0.72\unit{V}$, so
$I=0.72/0.30=\boxed{2.4\unit{A}}$, and the magnetic force is now
\emph{down} the incline, $BIL=1.2\cdot2.4\cdot0.30=0.864\unit{N}$. The
external force must overcome both:
\[
  F=mg\sin\theta+BIL=0.981+0.864=\boxed{1.85\unit{N}} .
\]
Power balance: the external power $Fv=1.85\cdot2.0=3.69\unit{W}$ splits
into $mg\sin\theta\,v=1.96\unit{W}$ of potential energy and
$I^2R=2.4^2\cdot0.30=1.73\unit{W}$ of Joule heat, and $1.96+1.73=3.69$.
\checkmark\ The magnetic force itself does no net work: it converts
mechanical energy into electrical energy at exactly the rate it is
dissipated.

\emph{iv.} With $\mathbf B$ vertical, the flux through the loop is
$BA_\perp$ with $A_\perp=A\cos\theta$, so
\[
  \mathcal{E}=\frac{d\Phi}{dt}=BLv\cos\theta . \tag{$\ast$}
\]
A current $I=BLv\cos\theta/R$ flows, and the magnetic force
$BIL$ is horizontal; its component along the rails is $BIL\cos\theta$, i.e.
\[
  F_{\parallel}=\frac{B^2L^2v\cos^2\theta}{R}.
\]
Setting $F_{\parallel}=mg\sin\theta$ gives
\[
  v_t=\frac{mgR\sin\theta}{B^2L^2\cos^2\theta}
  =\frac{v_{t,\perp}}{\cos^2\theta},
\]
so with $\theta=30^\circ$ the terminal speed is larger by $4/3$:
$\boxed{3.03\unit{m/s}}$. The extra factor of $1/\cos^2\theta$ has two
independent origins: only the component $B\cos\theta$ perpendicular to the
plane produces motional emf, and only the component $\cos\theta$ of the
magnetic force lies along the rails. Both must be compensated, hence the
square. (Check: with $\theta=0$ the expression reduces to the previous
case, as it must.)

\textbf{(b)} \emph{i.}
$\sin\theta_c=n_{cl}/n_c=1.45/1.46=0.99315$, so
$\theta_c=83.3^\circ$ from the normal, i.e.\ $6.7^\circ$ from the axis.
The numerical aperture is
\[
  \mathrm{NA}=\sqrt{n_c^2-n_{cl}^2}=\sqrt{2.916-2.9025}=0.171,
\]
so the acceptance angle in air is $\arcsin(\mathrm{NA})=9.8^\circ$. (The
numerical aperture is defined for the external medium; inside the core the
limiting ray makes $\arcsin(\mathrm{NA}/n_c)=6.7^\circ$ with the axis.)

\emph{ii.} At the limiting internal angle the geometric path in
$1.00\unit{m}$ of fibre is $1.00/\cos6.71^\circ=1.0069\unit{m}$, i.e.\
$0.69\unit{cm}$ extra. The delay relative to the axial ray is
\[
  \Delta t=\frac{\Delta L\,n_c}{c}
  =\frac{0.00690\cdot1.46}{3.00\cdot10^{8}}
  =\boxed{3.4\cdot10^{-11}\unit{s/m}},
\]
i.e.\ $34\unit{ps}$ per metre, or $34\unit{ns}$ per kilometre. Different
modes therefore arrive at slightly different times, so successive pulses
smear into each other (modal dispersion) and the bit rate is limited. In a
step-index fibre with a larger NA ($\approx0.3$) the spread is tens of
ns/km, which is severe. A graded-index core makes the longer paths travel through
lower-index glass near the cladding, where the light is faster, so the
transit times are approximately equalised.

\emph{iii.} $A/A_0=10^{-\alpha d/10}=10^{-0.25\cdot10/10}=10^{-0.25}
=\boxed{0.56}$, so about $44\%$ is lost over $10\unit{km}$. Silica
absorbs light by exciting electrons, which needs a photon energy above a
threshold --- in practice, in the ultraviolet. Violet photons ($h\nu
\approx2.5\unit{eV}$) are closer to that threshold and are scattered and
absorbed more strongly, while infrared photons ($\approx0.8\unit{eV}$) have
too little energy for electronic transitions and are absorbed only by weak
vibrational tails. That is why the attenuation minimum sits at
$1.3$--$1.55\unit{\mu m}$ and why long-haul fibres use infrared.

%======================================================================
\problem{10}{FY8}
\textbf{(a) i.} $E=hc/\lambda$ with $hc=1240\unit{eV\,nm}$:
\[
  E_1=\frac{1240}{589.0}=2.1050\unit{eV},\qquad
  E_2=\frac{1240}{589.6}=2.1029\unit{eV},
\]
so $\Delta E=2.1\unit{meV}$ and
\[
  \Delta f=\frac{\Delta E}{h}=\Delta E\cdot\frac{1\unit{eV}}{h}
  =\boxed{5.2\cdot10^{2}\unit{GHz}}.
\]
(Equivalently $\Delta f\approx f\,\Delta\lambda/\lambda
=5.09\cdot10^{14}\cdot(0.6/589)=5.2\cdot10^{11}\unit{Hz}$.)

\textbf{ii.} The grating constant is
$d=(1/600)\unit{mm}=1.667\unit{\mu m}$, and $d\sin\theta=k\lambda$:
\[
  \sin\theta_1=\frac{589.3\unit{nm}}{1.667\unit{\mu m}}=0.3536
  \Rightarrow\theta_1=20.7^\circ,
\]
\[
  \sin\theta_2=0.7071\Rightarrow\theta_2=45.0^\circ,\qquad
  \sin\theta_3=1.061>1\ \text{(no third order)} .
\]
So only orders $k=0,1,2$ exist. To resolve the doublet in first order,
$R=\lambda/\Delta\lambda=589.3/0.6=982=kN$ with $k=1$ requires
$N=\boxed{9.8\cdot10^{2}}$ illuminated lines, i.e.\ about
$1.6\unit{mm}$ of grating.

\textbf{iii.} Threshold wavelength
$\lambda_{\max}=hc/W=1240/2.28=\boxed{544\unit{nm}}$. At $500\unit{nm}$,
\[
  E_k^{\max}=\frac{1240}{500}-2.28=2.48-2.28=\boxed{0.20\unit{eV}},
\]
and the stopping potential is $\boxed{0.20\unit{V}}$. All three are
\emph{independent} of intensity; the intensity controls the \emph{number}
of photoelectrons per second (the photocurrent), not their maximum energy.
This is the content of Einstein's explanation: photon energy depends only
on frequency.

\textbf{(b) i.}
\[
  \Delta m=(2.014102+3.016049)-(4.002602+1.008665)
  =0.018884\unit{u},
\]
\[
  Q=0.018884\cdot931.494=\boxed{17.59\unit{MeV}} .
\]
Taking the deuteron as the consumed fuel,
\[
  \frac{Q}{m_D}
  =\frac{17.59\cdot1.6022\cdot10^{-13}\unit{J}}
         {2.014102\cdot1.6605\cdot10^{-27}\unit{kg}}
  =\boxed{8.4\cdot10^{13}\unit{J/kg}} .
\]
Fission of U-235 at $200\unit{MeV}$ per nucleus gives
$8.2\cdot10^{13}\unit{J/kg}$ of uranium, so \emph{per kilogram of fuel
consumed} the two are within a factor $10$ of each other --- the familiar
claim that fusion releases ``millions of times'' more energy is per kilogram
of the light \emph{reagents} only if one compares with chemical fuels, not
with fission. The decisive advantage of fusion is not the energy per kg but
the abundance of deuterium. Coal at $3.0\cdot10^{7}\unit{J/kg}$ is a factor
$2.8\cdot10^{6}$ worse: $1\unit{kg}$ of deuterium releases as much as about
$2.8\cdot10^{6}\unit{kg}$ of coal.

\textbf{ii.} With $P=2.5\cdot10^{9}\unit{W}$ and
$Q=17.59\unit{MeV}=2.818\cdot10^{-12}\unit{J}$,
\[
  \dot N=\frac{P}{Q}=\boxed{8.9\cdot10^{20}\unit{s^{-1}}} .
\]
In one year ($3.156\cdot10^{7}\unit{s}$) that consumes
$2.80\cdot10^{28}$ deuterons, i.e.\
\[
  m_D=2.80\cdot10^{28}\cdot2.014\unit{u}=\boxed{94\unit{kg}} .
\]
At $33\unit{g}$ of D per $\unit{m^3}$ of seawater this is about
$2.8\cdot10^{3}\unit{m^3}$ --- a swimming pool per GW-plant-year. Fuel
supply is emphatically \emph{not} the limiting problem; achieving and
sustaining the plasma conditions is.

\textbf{(c) i.} Here $t=12.0\unit{h}=0.500\unit{d}$, so
$t/T_{1/2}=0.500/8.02=0.06234$ and
\[
  A=A_0\,2^{-t/T_{1/2}}=20.0\cdot2^{-0.06234}
  =\boxed{19.2\unit{GBq}},
\]
a fall of only $4.2\%$. The number of nuclei falls by the same $4.2\%$,
which is why it ``does not visibly change'': $\lambda$ is tiny on this time
scale. With $\lambda=\ln2/T_{1/2}=1.0003\cdot10^{-6}\unit{s^{-1}}$ and
$N_0=A_0/\lambda$,
\[
  N_0=\frac{20.0\cdot10^{9}}{1.0003\cdot10^{-6}}
  =\boxed{2.0\cdot10^{16}} .
\]
That corresponds to only about $4\unit{\mu g}$ of ${}^{131}\mathrm{I}$:
the activity is high because the decay constant is large, not because much
material is present.

\textbf{ii.} Imaging needs radiation that \emph{leaves} the body.
${}^{131}\mathrm{I}$ is taken up by thyroid tissue, and its \emph{gamma}
rays ($364\unit{keV}$) penetrate tissue and reach an external camera; its
eight-day half-life matches the uptake biology while decaying away
afterwards. An $\alpha$ emitter would be stopped within a few
$\unit{\mu m}$ of the tissue that absorbed it --- no image, but a large
local dose. Alpha emitters are useful for targeted \emph{therapy}, not
imaging.

\textbf{(d)} Gamma: $H=0.20\cdot1=\boxed{0.20\unit{Sv}}$. Alpha:
$H=0.20\cdot20=\boxed{4.0\unit{Sv}}$, i.e.\ a potentially lethal equivalent
dose, twenty times the gamma figure for the same absorbed energy. The
factor $w_R=20$ already encodes the greater biological damage of densely
ionising tracks. What the grey-to-sievert conversion still ignores is the
\emph{spatial distribution} of the dose. A body-averaged $0.20\unit{Gy}$
from gammas is a fairly uniform exposure of all organs; the same average
delivered by alphas concentrated in the lungs means a much higher local
dose in a small, radio-sensitive tissue, and the relevant quantity becomes
the organ dose multiplied by a tissue weighting factor $w_T$ ($\approx
0.12$ for lung) rather than the whole-body average. In practice an alpha
dose is always local, so quoting a body-averaged number would understate
the local effect by orders of magnitude --- which is exactly why inhaled
alpha emitters such as radon daughters are so hazardous.

\end{document}
